\chapter{Applications, Evaluation, and Discussion}
\label{ch:evaluation}

% For theoretical work, use this chapter for examples, consequences, or
% limits. For empirical work, split it into protocol, results, and discussion.

Consider the graph in \cref{fig:example-graph}. The direct edges are
$a\to b$, $b\to c$, $c\to d$, and $a\to d$. The picture is drawn in TikZ, so
it stays sharp at any zoom level.

\begin{figure}[htbp]
  \centering
  \begin{tikzpicture}[
  vertex/.style={circle,draw=LMFINavy,thick,minimum size=8mm,inner sep=0pt},
  edge/.style={-{Stealth[length=2.4mm]},thick,LMFITeal},
  node distance=2.2cm
]
  \node[vertex] (a) {$a$};
  \node[vertex, right=of a] (b) {$b$};
  \node[vertex, right=of b] (c) {$c$};
  \node[vertex, below=1.0cm of c] (d) {$d$};

  \draw[edge] (a) -- (b);
  \draw[edge] (b) -- (c);
  \draw[edge] (c) -- (d);
  \draw[edge] (a) to[bend right=18] (d);
\end{tikzpicture}

  \caption{A four-vertex directed graph used to evaluate reachability.}
  \label{fig:example-graph}
\end{figure}

Applying \cref{alg:warshall} gives the reflexive reachability matrix in
\cref{tab:reachability}. For instance, $a$ reaches every vertex, whereas $d$
reaches only itself. The entry in row $b$ and column $d$ records the path
$b\to c\to d$, even though there is no direct edge $b\to d$.

\begin{table}[htbp]
  \centering
  \caption{Reflexive reachability matrix for \cref{fig:example-graph}.}
  \label{tab:reachability}
  \begin{tabular}{c *{4}{>{\centering\arraybackslash}p{1.2cm}}}
    \toprule
    From $\backslash$ To & $a$ & $b$ & $c$ & $d$ \\
    \midrule
    $a$ & 1 & 1 & 1 & 1 \\
    $b$ & 0 & 1 & 1 & 1 \\
    $c$ & 0 & 0 & 1 & 1 \\
    $d$ & 0 & 0 & 0 & 1 \\
    \bottomrule
  \end{tabular}
\end{table}

This example checks internal consistency rather than performance. A genuine
experimental chapter should state the data, software version, metric, and
uncertainty, and discuss negative cases: reachability alone loses path lengths.
